\documentclass{jsarticle}
\usepackage[dvipdfmx]{graphicx}
\usepackage[LY1]{fontenc}
\usepackage{amsmath,amsfonts}
\usepackage{amssymb}
\usepackage[mtbold,mtpluscal]{mathtime}
\usepackage[scaled]{helvet}
\usepackage{bm}
\begin{document}
\title{}
\author{}
\date{}
\maketitle


 \newcommand{\letter}{%
	 \begin{picture}(10,10)
		 \put(0,0){\line(1,1){10}}
		 \put(5,5){\line(-1,1){10}}
		 \put(10,10){\line(0,-0){10}}
	 \end{picture} %
	 }
	 \letter


 \newcommand{\aletter}{%
	 \begin{picture}(10,10)
		 \put(0,0){\line(1,1){10}}
		 \put(5,5){\line(-1,1){10}}
		 \put(10,10){\line(-5,0){10}}
	 \end{picture} %
	 }
	 \aletter
 \newcommand{\eletter}{%
	 \begin{picture}(10,10)
		 \put(0,0){\line(1,1){10}}
		 \put(10,10){\line(1,-1){10}}
	 \end{picture} %
	 }
	 \eletter
 \newcommand{\siletter}{%
	 \begin{picture}(10,10)
		 \put(0,0){\line(1,1){10}}
		 \put(0,0){\line(0,10){10}}
		 \put(0,10){\line(20,0){10}}
		 \put(0,0){\line(0,20){10}}
		 \put(10,0){\line(-20,0){10}}
		 \put(10,10){\line(1,-1){10}}
	 \end{picture} %
	 }
	 \siletter

 \newcommand{\sletter}{%
	 \begin{picture}(10,10)
		 \put(0,0){\line(0,10){10}}
		 \put(0,10){\line(20,0){10}}
		 \put(0,0){\line(0,20){10}}
		 \put(10,0){\line(-20,0){10}}
		 \put(10,10){\line(0,-10){10}}
	 \end{picture} %
	 }
	 \sletter
 \newcommand{\tletter}{%
	 \begin{picture}(10,10)
		 \put(0,0){\line(0,10){10}}
		 \put(0,10){\line(1,-1){10}}
		 \put(10,0){\line(-1,1){10}}
		 \put(0,10){\line(20,0){10}}
		 \put(0,0){\line(0,20){10}}
		 \put(10,0){\line(-20,0){10}}
		 \put(10,10){\line(0,-10){10}}
	 \end{picture} %
	 }
	 \tletter


 \newcommand{\ateletter}{%
	 \begin{picture}(10,10)
		 \put(0,0){\line(0,10){10}}
		 \put(0,10){\line(1,-1){10}}
		 \put(10,10){\line(-1,-1){10}}
		 \put(0,10){\line(20,0){10}}
		 \put(0,0){\line(0,20){10}}
		 \put(10,0){\line(-20,0){10}}
		 \put(10,10){\line(0,-10){10}}
	 \end{picture} %
	 }
	 \ateletter

 \newcommand{\vletter}{%
	 \begin{picture}(20,20)
		 \put(0,0){\line(0,20){20}}
		 \put(0,20){\line(1,-1){10}}
		 \put(20,20){\line(-1,-1){10}}
		 \put(0,20){\line(20,0){20}}
		 \put(0,0){\line(0,20){20}}
		 \put(20,0){\line(-20,0){20}}
		 \put(20,20){\line(0,-20){20}}
	 \end{picture} %
	 }
	 \vletter

 \newcommand{\variiint}{%
	 \begin{picture}(30,30)
		 \put(0,0){$ \iiint $}
		 \put(0,5){\line(20,0){20}}
		 
	 \end{picture} %
	 }
	 \variiint







	 




\end{document}
